% !TEX root = ../../tree_notes.tex
% =====================================================================
%  9.2  Monotonic Stack Applications            sec:c26-monostack
%  Exports: def:c26-monostack  thm:c26-hist  fig:c26-histfig
%           lst:c26-hist  lst:c26-deque  ex:c26-hist-run
%           prob:c26-hist84  lst:c26-hist84-soln  (practices: LC 85, 42, 239, 739)
% =====================================================================
\section{Monotonic Stack Applications}\label{sec:c26-monostack}

\begin{previouslybox}
\begin{enumerate}[nosep, leftmargin=1.6em]
  \item What two orders define a Cartesian tree, and where is its root? (\defref{c26-cartesian}; the
        minimum, by the heap rule)
  \item State \lemref{c26-span} in one sentence: node $i$'s subtree is the window over which $a[i]$
        is the \emph{what}, bounded by the nearest \emph{what} on each side?
  \item The build \coderef{c26-stackbuild} keeps the stack's values \emph{increasing}. When does index
        $i$ get pushed, and which node does it pop past? (past the run of larger values)
  \item Why is a whole monotonic sweep $\bigO{n}$ even with a nested \code{while}-pop loop? (each index
        pushed once, popped at most once -- the amortised-stack argument)
\end{enumerate}
\end{previouslybox}

\begin{topicintro}
Strip the Cartesian tree back to its mechanics and you are left with one idea so small it fits in a
\code{while} loop: keep a stack whose values are \emph{monotone}, and whenever a new element breaks the
monotonicity, pop the run it breaks -- and \emph{at the instant of popping, the popped element's answer
is finally known}. That is the \term{monotonic stack}, and it is the workhorse of this whole chapter.
The popped element has just met its nearest greater (or smaller) neighbour: the exact bound
\lemref{c26-span} needed.

This section makes the stack a technique in its own right -- you no longer need to name the tree -- and
runs the entire application family on it: the largest rectangle in a histogram (the solved anchor,
LeetCode 84), the next-greater family (daily temperatures), sliding-window extremes (a monotonic
\emph{deque}, LeetCode 239), and trapping rain water. The unifying sentence, which you should memorise:
\textbf{a bar's rectangle is decided the moment it is popped, because that is when the shortest bar on
either side of its dominance window are both known}.
\textbf{Prerequisites.} \secref{c26-cartesian} (the span lemma is the reason the stack works) and a
\code{stack}/\code{deque} container.

\textbf{After this section you can}
\begin{itemize}[nosep]
  \item decide the pop-comparison ($>$ vs $\ge$, $<$ vs $\le$) that matches the bounds a problem needs; \hfill\lv{3}
  \item solve largest-rectangle-in-histogram cold, and explain why the width is read at pop time; \hfill\lv{4}--\lv{5}
  \item convert any "next greater / smaller element" problem to one decreasing or increasing stack pass; \hfill\lv{3}
  \item keep a sliding-window extreme with a monotonic deque at amortised $\bigO{1}$ per position; \hfill\lv{4}
\end{itemize}
\end{topicintro}

\begin{recognitionbox}[Pattern card -- nearest bigger/smaller neighbour]
\reckey{each element's fate is sealed by the \emph{first} element to its left or right that is
\emph{larger} (or \emph{smaller}); or a \emph{window extreme} that must be reported as the window slides.}
\recuse{
  \item "largest rectangle under a skyline", "maximal rectangle of 1s in a grid" -- pop when the bar
        stops being the limiting height (\probref{c26-hist84}, LC 85);
  \item "next greater element", "days until a warmer temperature" -- one decreasing stack;
  \item "max (min) in every window of size $k$" -- monotonic \emph{deque} (\coderef{c26-deque});
  \item "water trapped between bars" -- stack of walls, or two-pointer prefix maxima (LC 42).
}
\recnotuse{
  \item you need arbitrary-range extremes over a \emph{dynamic} array -- segment tree (\chref{c13}),
        the stack cannot forget its pops;
  \item the window extreme is the \emph{median} -- a monotonic deque tracks only the extreme, not the
        middle; that is a multiset, Chapter~8 (\chref{c06});
  \item "the $k$-th smallest overall" -- a heap (\chref{c07}), not a stack.
}
One-line test: \emph{does the answer for an element depend only on the nearest element that beats it?
Then a monotonic stack finds every such answer in one pass.}
\end{recognitionbox}

\subsection{Intuition}

Imagine the histogram $H=(2,1,5,6,2,3)$ of \figref{c26-hist} as a city skyline and drop a vertical
waterline at some height $h$. At height $h$, the tallest rectangle of height $h$ is limited by the
\emph{nearest bar shorter than $h$} on the left and the nearest bar shorter than $h$ on the right --
everything between them is at least $h$. So for the bar at position $b$, the rectangle of height
$H[b]$ that uses bar $b$ as its \emph{limiting} height spans exactly those two shorters' inner window:
this is \lemref{c26-span} again, now read as ``how far can a bar of height $H[b]$ extend?''

The monotonic stack computes all these windows simultaneously. Walk the bars left to right keeping a
stack of positions whose heights are \emph{strictly increasing} -- a rising staircase. When a new bar
$i$ is \emph{lower} than the stack top, the top bar has just met its right-boundary shorter neighbour:
pop it, and its left-boundary shorter neighbour is now the (new) stack top. Between those two bounds,
its rectangle is fully determined, so evaluate its area \emph{right now}. Bars that never get popped
until the end are popped by a zero-height sentinel appended to the array. One pass, every rectangle --
this is the same loop as \coderef{c26-stackbuild}, with the "compute the span's area" step surfacing.

\begin{figure}[H]
\centering
\begin{tikzpicture}[diagram, x=11mm, y=6mm]
  \draw[thick, darkgray] (0,0) -- (6.4,0);
  \foreach \i/\h in {0/2,1/1,2/5,3/6,4/2,5/3} {
    \fill[primaryblue!22, draw=primaryblue, thick]
      (\i,0) rectangle (\i+1,\h);
    \node[font=\scriptsize, text=darkgray] at (\i+0.5,-0.45) {\i};
    \node[font=\scriptsize\bfseries, text=deepblue] at (\i+0.5,\h+0.3) {\h};}
  % the limiting window for the height-5 bar: bars 2 and 3, blocked by bar1(=1) and bar4(=2)
  \draw[ultra thick, warningred] (2,0) rectangle (4,5);
  \draw[warningred, dashed] (1,5) -- (4,5);
  \node[font=\scriptsize\bfseries, text=warningred, align=center] at (3,2.5) {$h{=}5$\\ width $=2$};
  \node[annot, align=left, font=\scriptsize] at (7.0,2.5)
    {bar $2$ pops at bar $4$:\\ left bound $=$ bar $1$ ($h{=}1$)\\ right bound $=$ bar $4$ ($h{=}2$)\\
     width $4-1-1=2$, area $10$};
\end{tikzpicture}
\caption{Why the width is decided \emph{at pop time}. Bar $2$ (height $5$) sits on the increasing stack
until bar $4$ (height $2$) is lower and pops it. At that instant its left-boundary shorter bar is the
new stack top, bar $1$ (height $1$), and its right-boundary shorter bar is $i=4$; the rectangle of
height $5$ spans bars $2$--$3$, width $4-1-1=2$, area $10$. That is \thmref{c26-hist} and the answer for
$H$.}\label{fig:c26-histfig}
\end{figure}

\subsection{Definitions and Notation}

\begin{definition}{Monotonic stack}{c26-monostack}
A \term{monotonic stack} is a stack of indices whose stored values are kept \emph{non-decreasing} (an
\emph{increasing} stack, used to find nearest \emph{smaller} elements) or \emph{non-increasing} (a
\emph{decreasing} stack, for nearest \emph{greater} elements). When a new index $i$ arrives, pop from
the top while the invariant would be violated; each popped index has just discovered its nearest
beating neighbour ($i$ to its right), and the new top is its neighbour to the left. Every index enters
and leaves the stack once, so the total work is $\bigO{n}$.
\end{definition}

\noindent\textbf{In plain words.} The stack is a queue of elements still ``waiting'' to meet someone
bigger/smaller than them. The moment such someone appears, all the elements it beats are answered at
once and leave. Nothing is ever looked at twice.

\begin{theorem}{A popped bar's rectangle is finalised at pop}{c26-hist}
Run an increasing stack (pop while $H[\mathit{top}] \ge$ the incoming height) over $H$ with a
zero-height sentinel appended. When bar $b$ is popped at index $i$, let $L$ be the new stack top
($L=-1$ if the stack became empty). Then $i$ is the nearest index to the right of $b$ with a height
\emph{strictly less} than $H[b]$ reached, $L$ is the nearest to the left, and the largest rectangle of
height $H[b]$ has width $i - L - 1$ and area $H[b]\cdot(i-L-1)$. Taking this over every pop yields the
maximum rectangle in the histogram.
\end{theorem}
\begin{proofbox}
Bars are pushed in increasing height order, so while $b$ stayed on the stack every bar pushed after it
and popped before $i$ had height $\ge H[b]$ and lies inside $b$'s window. The pop triggers at $i$ exactly
when $H[i] < H[b]$ (or the sentinel), so $i$ is the first shorter bar to the right; after popping $b$,
the stack top $L$ is the last bar with $H[L] < H[b]$ to the left (everything between $L$ and $b$ has
height $\ge H[b]$ or was already popped as taller). Hence the maximal contiguous run of bars of height
$\ge H[b]$ that includes $b$ is $[L+1, i-1]$, of width $i-L-1$; the rectangle of height $H[b]$ is
attained over exactly that run. \emph{Every} candidate optimal rectangle has some height equal to one of
the bars (raise it until it hits a bar top), say bar $b$, and is at most the width just computed, so the
best over all pops equals the true optimum. \end{proofbox}

\begin{insightbox}
\thmref{c26-hist} is \lemref{c26-span} wearing a hat. ``The subtree of $b$ is its dominance window'' and
``the rectangle of bar $b$ is bounded by the nearest shorter bars'' are the \emph{same statement} -- a
shorter bar is exactly a smaller value that becomes an ancestor. If you internalised 9.1, the histogram
is already solved; this section just shows the stack \emph{without} drawing any tree, which is how you
will actually write it in contest code.
\end{insightbox}

\subsection{The algorithm: one pass over the bars}

The code is the sentinel-augmented increasing stack. Keep positions, not heights, so you can recover
the width from the index gap; pop on $\ge$ so a taller-or-equal bar is forced to settle its rectangle
before an equal bar can widen it.

\begin{lstlisting}[caption={Largest rectangle in a histogram -- one increasing stack with a zero sentinel}, label={lst:c26-hist}]
long long largestRectangle(const vector<long long>& H) {
    int n = (int)H.size();
    vector<int> st;                       // indices, heights strictly rising
    long long best = 0;
    for (int i = 0; i <= n; ++i) {        // i == n is the 0-height sentinel
        long long cur = (i == n ? 0 : H[i]);
        while (!st.empty() && H[st.back()] >= cur) {
            long long h = H[st.back()]; st.pop_back();
            int L = st.empty() ? -1 : st.back();   // new top = left shorter bound
            best = max(best, h * (i - L - 1));      // width decided at the pop
        }
        st.push_back(i);
    }
    return best;
}
\end{lstlisting}
\begin{complexitybox}
Time $\bigO{n}$ -- each index pushed once and popped at most once. Memory $\bigO{n}$ for the stack.
Verdict: the linear skyline algorithm; the naive $\bigO{n^2}$ (grow a window from each bar, track its
running minimum) is the same answer an order slower and \tle s for $n\sim10^5$.
\end{complexitybox}

\subsection{Worked Example}

\begin{example}{The stack on $H=(2,1,5,6,2,3)$}{c26-hist-run}
Read the pop events; $\mathit{st}$ holds indices and the sentinel is at $i=6$ with height $0$.\\
$i{=}0$ push; $i{=}1$ (height $1$) pops bar $0$ ($h=2$, $L=-1$, width $2$? no -- $1-(-1)-1=1$, area $2$),
then pushes $1$. $i{=}2,3$ push (heights $5,6$ rising). At $i{=}4$ (height $2$): pop bar $3$ ($h=6$,
$L=2$, width $4-2-1=1$, area $6$); pop bar $2$ ($h=5$, $L=1$, width $4-1-1=2$, area $10$); push $4$.
$i{=}5$ (height $3$) pushes. Sentinel $i{=}6$: pop bar $5$ ($h=3$, $L=4$, width $1$, area $3$); pop bar
$4$ ($h=2$, $L=1$, width $6-1-1=4$, area $8$); pop bar $1$ ($h=1$, $L=-1$, width $6$, area $6$).
Maximum over the areas $\{2,6,10,3,8,6\}$ is $\boxed{10}$ -- the height-$5$ rectangle over bars $2,3$. \checkmark
\end{example}
\begin{tracebox}[Pop ledger -- each bar settles when a shorter one arrives]
bar 0 (h2): popped by bar1, width 1, area 2.\\
bar 1 (h1): popped by sentinel, width 6, area 6 (spans everything -- height 1 reaches all).\\
bar 2 (h5): popped by bar4, width 2, area 10.\\
bar 3 (h6): popped by bar4, width 1, area 6.\\
bar 4 (h2): popped by sentinel, width 4, area 8 (bars 2..5 minus the height-1 blocker at bar1).\\
bar 5 (h3): popped by sentinel, width 1, area 3.  Best area $=10$.
\end{tracebox}

\subsection{The wider family}

The exact same loop, with the comparison flipped and the ``what do you compute at pop'' changed, mops up
a shelf of problems. Keep the four dialects straight:

\begin{itemize}[nosep]
  \item \textbf{Next greater element} (\emph{decreasing} stack). Pop while the top's value is $<$ the
        incoming one; the incoming index is that element's next greater. Daily temperatures
        (LC 739) is literally ``distance to next greater''. \hfill\lv{3}
  \item \textbf{Previous smaller} (\emph{increasing} stack). The pop you already do in the histogram
        gives the nearest shorter bar to the left for free -- this is $L$ in \thmref{c26-hist}. \hfill\lv{3}
  \item \textbf{Sliding-window maximum} (\emph{monotonic deque}, LC 239). Same decreasing stack, but you
        also drop the \emph{front} once it leaves the window of size $k$; a deque pops both ends. \hfill\lv{4}
  \item \textbf{Trapping rain water} (LC 42). A decreasing stack of walls catches valleys: pop a floor,
        the current bar and the new top are its bounding walls, fill the basin. (Or two-pointer prefix
        maxima -- same idea, see Practice.) \hfill\lv{4}--\lv{5}
\end{itemize}

\begin{lstlisting}[caption={Sliding-window maximum with a monotonic deque (LeetCode 239) -- drop smaller from the back, out-of-window from the front}, label={lst:c26-deque}]
vector<int> maxSlidingWindow(vector<int>& nums, int k) {
    deque<int> dq;                       // indices, values non-increasing, front = current max
    vector<int> out;
    for (int i = 0; i < (int)nums.size(); ++i) {
        while (!dq.empty() && nums[dq.back()] <= nums[i]) dq.pop_back();  // dead: i beats them
        dq.push_back(i);
        if (dq.front() <= i - k) dq.pop_front();                          // left the window
        if (i >= k - 1) out.push_back(nums[dq.front()]);                  // window is full
    }
    return out;
}
\end{lstlisting}

\subsection{Common Pitfalls}\label{sec:c26-monostack-pitfalls}

\begin{warningbox}[Trap 1 -- $>$ versus $\ge$ silently miscounts widths]
Popping while $H[\text{top}] > \text{cur}$ (strict) leaves \emph{equal} heights stacked, and the left
bound $L$ for a popped bar is then an equal bar -- so the width is under-counted until a later equal bar
pays it back, and the \emph{maximum} still comes out right. But the moment you \emph{sum} areas or count
ranges (LC 907, the Summit) the same off-by-relationship double- or under-counts. Pick one convention per
side deliberately -- here $\ge$ on the pop keeps the max correct and each rectangle's true width is
recovered when its \emph{last} equal-height bar pops.
\end{warningbox}
\begin{lstlisting}[caption={\textcolor{warningred}{WRONG}: no sentinel -- tall bars never get their right bound}, language=C++, basicstyle=\ttfamily\small]
for (int i = 0; i < n; ++i) { ... }      // stack may still hold bars when the loop ends
// the tallest bars (right of every shorter bar) are never popped -> their rectangles uncomputed
\end{lstlisting}
\begin{lstlisting}[caption={\textcolor{successgreen}{RIGHT}: run to i==n with cur==0 to flush the stack}, language=C++, basicstyle=\ttfamily\small]
for (int i = 0; i <= n; ++i) {           // i == n appends a height-0 sentinel
    long long cur = (i == n ? 0 : H[i]); // forces every remaining bar to settle its width
    ...
}
\end{lstlisting}

\begin{warningbox}[Trap 2 -- storing values instead of indices]
The width is an \emph{index} difference ($i - L - 1$), so the stack must hold positions, not heights. If
you push heights you lose the ability to measure how far the rectangle reaches -- the whole point of
\thmref{c26-hist}. Keep indices and read $H[\text{index}]$ when you need the height.
\end{warningbox}

\begin{warningbox}[Trap 3 -- a deque for the wrong window statistic]
A monotonic deque answers sliding-window \emph{max/min} in $\bigO{1}$ amortised because you may discard
any element that can never again be the extreme. You cannot do that for the sliding \emph{median} or
\emph{sum} -- smaller elements still matter for the median -- so those need an ordered multiset or two
heaps (Chapter~8, \chref{c06}; Chapter~10, \chref{c07}). Do not reach for the deque out of habit.
\end{warningbox}

\subsection{Problem Lab}

The section's anchor is the canonical LeetCode 84 -- the histogram running example. It is chosen because
it is \thmref{c26-hist} with the proof already done: one increasing stack, a sentinel, and a multiply.

% ---------------------------------------------------------------------------
\solvedproblem{c26-hist84}{Largest Rectangle in Histogram}{LeetCode 84}
  {https://leetcode.com/problems/largest-rectangle-in-histogram/}{LC Hard · L5}
  {monotonic (increasing) stack, sentinel flush, width read at pop time}

\begin{problemstatement}[Largest Rectangle in Histogram]
Given an array of $n$ non-negative integers $\mathit{heights}$ representing a histogram where the width
of each bar is $1$, return the area of the \textbf{largest rectangle} that can be formed within the
histogram. $1 \le n \le 10^{5}$ and $0 \le \mathit{heights}[i] \le 10^{4}$.
\end{problemstatement}

\begin{attemptgate}[Attempt it yourself first -- \lv{5}, give yourself 30 minutes]
The answer is a rectangle whose height equals \emph{some} bar and whose width is how far that bar's
height can extend left and right before a shorter bar blocks it.
\peek{Fix a bar $b$. What bounds the widest rectangle of height $H[b]$? (Nearest bar shorter than $H[b]$
on each side -- \thmref{c26-hist}.)}
\nudge{Sweep left to right with an increasing stack of indices. When a shorter bar arrives it pops the
taller ones, and for each popped bar you now know both bounds -- compute its rectangle immediately.
Append a $0$-height bar at the end to flush whatever is left.}
\fullhint{On pop of bar $x$ at incoming index $i$: $L$ is the new stack top ($-1$ if empty), width $=
i - L - 1$, area $= H[x]\cdot\text{width}$. Take the max over all pops. $\bigO{n}$. For the sample the
height-$5$ bar pops giving width $2$, area $10$.}
\end{attemptgate}

\solstep{Understanding the Problem}
A rectangle in the histogram is a contiguous set of bars lifted to a common height; that height cannot
exceed the shortest bar in the set. So an optimal rectangle has height exactly $H[b]$ for some bar $b$
(its limiting bar) and width equal to the maximal run around $b$ of bars all $\ge H[b]$. The task is thus:
for every bar, find its left and right \emph{nearest strictly shorter} bars, and evaluate
$H[b]\cdot\text{width}$. Doing that naively per bar is $\bigO{n}$ each, $\bigO{n^2}$ overall; the
monotonic stack finds every bar's two bounds in one shared pass.

\solstep{Key Observations}
\begin{observation}
On an increasing stack of indices, when bar $x$ is popped by an incoming bar $i$, $i$ is $x$'s nearest
shorter bar on the \emph{right} and the new stack top $L$ is $x$'s nearest shorter bar on the
\emph{left}. Both bounds become known at the single moment of the pop (\thmref{c26-hist}).
\end{observation}
\begin{observation}
The rectangle of height $H[x]$ therefore spans indices $[L+1, i-1]$, width $i - L - 1$, area
$H[x]\cdot(i-L-1)$. No bar between $L$ and $i$ is shorter than $H[x]$ (they were pushed after $x$ or popped as taller).
\end{observation}
\begin{observation}
Appending a $0$-height sentinel guarantees every bar is eventually popped, so no candidate rectangle is
missed even for a strictly-increasing skyline.
\end{observation}

\solstep{Think It Through With Me}
\begin{thinkbox}
\deadend{For each bar, walk left and right counting bars of height $\ge$ it.}
{$\bigO{n^2}$: on a flat histogram of height $1$ and $n=10^5$, every bar scans all $n$, giving $10^{10}$
-- \tle. It also recomputes the same bounds over and over.}
\deadend{Divide and conquer on the minimum (Cartesian tree + RMQ).}
{Correct and $\bigO{n\log n}$ (build the tree, recurse on min), and it \emph{is} the same structure --
but the constant and the recursion are unnecessary when one flat stack pass is $\bigO{n}$.}
\aha{One increasing stack with a sentinel: a bar is popped exactly when its right bound appears, its left
bound is the new top, and the area is a multiply. $\bigO{n}$, in-place, no tree built.}
\end{thinkbox}

\solstep{Solution}
Traverse $i$ from $0$ to $n$ inclusive. Let $\mathit{cur}$ be $H[i]$ (or $0$ at the $i=n$ sentinel). While
the stack is non-empty and $H[\text{top}] \ge \mathit{cur}$: pop index $x$; set $L = \text{new top}$ or
$-1$ if empty; $\mathit{best} = \max(\mathit{best}, H[x]\cdot(i - L - 1))$. Then push $i$. Return
$\mathit{best}$ (as a \code{long long} if the product can exceed $2^{31}$).

\solstep{Proof of Correctness}
\begin{proofbox}[The stack attains the true maximum]
(\emph{Each pop evaluates a valid rectangle}) By \thmref{c26-hist}, at the pop of $x$ the bounds $L$ and
$i$ are the nearest strictly-shorter bars on each side, so $[L+1, i-1]$ is the widest run of bars of
height $\ge H[x]$ containing $x$; the area $H[x]\cdot(i-L-1)$ is a genuinely realisable rectangle and the
running maximum never exceeds the true optimum. (\emph{Nothing is missed}) Any optimal rectangle has
height equal to the shortest bar $x$ it covers; by the sentinel argument $x$ is popped at some index
which is precisely its right shorter bound, and at that pop its left bound $L$ is exactly its left
shorter bound, so the algorithm evaluates the optimum's own area. Hence $\mathit{best}$ equals the true
maximum. \end{proofbox}

\solstep{Implementation}
\begin{lstlisting}[caption={Largest Rectangle in Histogram -- the sentinel monotonic stack (LeetCode 84)}, label={lst:c26-hist84-soln}]
class Solution {
public:
    int largestRectangleArea(vector<int>& heights) {
        int n = (int)heights.size();
        vector<int> st;                          // increasing stack of indices
        long long best = 0;
        for (int i = 0; i <= n; ++i) {           // i == n is the 0-height sentinel
            int cur = (i == n ? 0 : heights[i]);
            while (!st.empty() && heights[st.back()] >= cur) {
                long long h = heights[st.back()]; st.pop_back();
                int L = st.empty() ? -1 : st.back();
                best = max(best, h * (i - L - 1));
            }
            st.push_back(i);
        }
        return (int)best;
    }
};
\end{lstlisting}

\solstep{Trace}
\begin{tracebox}[Official sample -- heights \texttt{[2,1,5,6,2,3]} (answer 10)]
See \exref{c26-hist-run}. The decisive pop is bar index $2$ (height $5$) when the shorter bar index $4$
(height $2$) arrives: $L=1$ (the height-$1$ bar), width $4-1-1=2$, area $5\cdot2=10$. Every other
settling area is $\le 10$, so the answer is $10$ -- the red rectangle in \figref{c26-histfig}.
\end{tracebox}

\solstep{Complexity}
\begin{complexitybox}
Time $\bigO{n}$: each index pushed once, popped at most once, so the inner \code{while} totals $\le n$
iterations across the whole sweep. Memory $\bigO{n}$ for the stack. Verdict: \ac -- the linear histogram
algorithm. Its $\bigO{n^2}$ sibling (grow a window per bar) fails at $n=10^5$.
\end{complexitybox}

\solstep{Pitfalls and Variants}
\begin{warningbox}
Ways to fall: (i) forgetting the $i=n$ sentinel -- tall trailing bars never settle their width; (ii)
pushing \emph{heights} not \emph{indices} -- no way to measure the width (Trap~2); (iii) \code{int}
overflow in $h\cdot\text{width}$ when the grid is wide -- accumulate in \code{long long}. \emph{Variant
(LC 85, Maximal Rectangle):} treat each matrix row as a histogram of consecutive 1s above it and call
this routine per row -- the largest all-1 rectangle is the largest rectangle over those stacked
histograms.
\end{warningbox}

\begin{insightbox}[Key takeaways]
\begin{enumerate}[nosep, leftmargin=1.6em]
  \item \textbf{Pop time $=$ answer time.} An element's nearest beating neighbour is known the instant it
        is popped; anything computed from that neighbour (width, next-greater distance, basin) is a
        multiply or subtract away (\thmref{c26-hist}).
  \item \textbf{One loop, four dialects.} Increasing vs decreasing stack and what you do at the pop
        decide whether you get histogram, next-greater, sliding max, or rain water.
  \item \textbf{A deque is a stack that also pops the front.} That extra move is exactly the
        ``out-of-window'' eviction that turns the stack into a sliding-window extreme
        (\coderef{c26-deque}).
\end{enumerate}
\end{insightbox}

\begin{hardeningbox}
\begin{itemize}[nosep, leftmargin=1.4em]
  \item \textbf{Next smaller to both sides in one pass} -- the pop gives right bound, the new top the left
        bound; return both arrays and you have \lemref{c26-span}'s $L,R$ for a sum problem.
  \item \textbf{Maximum of minimums of every window size} (a classic) -- run the histogram span, then
        bucket each bar's area by its width and sweep from the right keeping the max.
  \item \textbf{Stock-span / online variants} -- feed the decreasing stack as data arrives; the same pop
        counts how far back the current element dominates.
\end{itemize}
\end{hardeningbox}

\subsection{Practice Problems}

\begin{stuckbox}
\begin{itemize}[nosep, leftmargin=1.4em]
  \item Did you append the sentinel (or equivalently flush the stack after the loop)?
  \item Did the stack hold \emph{indices} so widths are index differences?
  \item For next-greater: did you choose a \emph{decreasing} stack so the arriving bigger element answers
        the popped one?
\end{itemize}
\end{stuckbox}

\begin{practiceset}
  \practice{Maximal Rectangle}{LeetCode 85}
    {https://leetcode.com/problems/maximal-rectangle/}{LC Hard · L4}
    {Given an $n\times m$ binary matrix, return the area of the largest rectangle containing only 1s.
    Turn each row into a histogram: $\mathit{height}[j]$ is the number of consecutive 1s in column $j$
    up to and including this row (extend if this row's cell is 1, else reset to 0). Then run the solved
    \probref{c26-hist84} on that histogram; the best over all rows is the answer. The key realisation --
    a rectangle of 1s ending at row $r$ is exactly a histogram rectangle in $r$'s stacked skyline.}

  \practice{Trapping Rain Water}{LeetCode 42}
    {https://leetcode.com/problems/trapping-rain-water/}{LC Hard · L4}
    {Given bar heights, how much water is trapped after raining. Two linear routes: the \emph{two-pointer}
    method keeps a left-max and right-max and fills each cell by $\min(\text{leftMax},\text{rightMax})-\text{height}$;
    or the \emph{monotonic decreasing stack}: when a taller bar pops a lower one, the popped slot is a
    basin floor bounded left by the new stack top and right by the arriving bar -- fill width $\times$
    (min of the two walls $-$ floor). Both are $\bigO{n}$; the stack version is the same pop-decides-answer
    machinery as \thmref{c26-hist}.}

  \practice{Sliding Window Maximum}{LeetCode 239}
    {https://leetcode.com/problems/sliding-window-maximum/}{LC Hard · L4}
    {Return the maximum in each window of size $k$ as it slides across the array. Use the monotonic
    deque of \coderef{c26-deque}: keep indices with non-increasing values; push new $i$ after popping
    smaller-or-equal from the back, evict the front once its index $\le i-k$, and read the front once the
    window is full. Amortised $\bigO{1}$ per element -- a heap with lazy deletion is the $\bigO{n\log n}$
    fallback but the deque is optimal.}

  \practice{Daily Temperatures}{LeetCode 739}
    {https://leetcode.com/problems/daily-temperatures/}{LC Medium · L3}
    {For each day, how many days until a strictly warmer temperature; $0$ if none. This is next-greater
    distance: keep a \emph{decreasing} stack of indices; when a warmer day arrives, pop every cooler index
    and answer it as $i - \text{popped}$. Push each index once, pop once -- $\bigO{n}$. The cleanest
    possible ``the pop is the answer'' warm-up before LC 84.}
\end{practiceset}

\begin{recallbox}[10-second refresher -- 9.2 Monotonic Stack Applications]
A monotonic stack keeps its values ordered so that an element is popped the moment it meets the nearest
neighbour that beats it -- and the answer for that element is finalised \emph{at the pop}
(\defref{c26-monostack}, \thmref{c26-hist}). Increasing stack $\to$ nearest smaller (histogram,
\coderef{c26-hist}); decreasing stack $\to$ nearest greater (daily temperatures). Append a sentinel to
flush the stack. Add a front-eviction and the stack becomes a monotonic \emph{deque} for sliding-window
extremes (\coderef{c26-deque}, LC 239). Everything is amortised $\bigO{n}$. LC 84 is the anchor; LC 85,
42, 239, 739 are the dialects.
\selftest{
  \item Why does the width of a popped bar's rectangle become known exactly at its pop?
  \item Increasing or decreasing stack for "next greater element"? Why?
  \item What does the $\bigO{1}$ sliding-window max deque pop, and from which ends?
  \item Give the rain-water pop recipe: for a popped basin floor, what are the two walls?
  \item Why does forgetting the sentinel under-report areas on a rising skyline?
}
\end{recallbox}
